Selecting source-data pools for a target task can require training and validating a predictor for many combinations of pools. Frozen features allow an earlier comparison, but retaining promising candidates is useful only if the final choice and the work saved justify screening. We evaluate a two-stage protocol that uses unlabeled source and target features to form a shortlist, then trains predictors on source labels and selects one using target validation. We compare target-weighted experimental design (A-opt) with second-moment matching (MMD) and target-blind diversity. The analysis explains the design score under a shared linear model and separates candidate omission from validation error. On equal-sample-cost DomainNet combinations, evaluating 50 A-opt or MMD candidates gives lower mean test Brier than evaluating 227 random candidates. Moment matching remains competitive across exploratory projection and readout changes. Domain-level matching also performs well after fixed-parameter repartitioning, while within-domain gains depend on construction and loss metric. Direct-selection gains on PACS and Office-Home are small, with limited improvement headroom. Independent timing shows savings with cached features. Adding measured feature-extraction costs preserves the logistic-regression saving but reverses it for ridge regression. These results assess screening through final selection quality and computation cost, rather than candidate retention alone.
